% ID: SEL-D0038-OA-E03
% ESTADO: validada. S039, cotejada con examen y criterios.
Sean $L$ llover y $V$ hacer viento. Para un día elegido al
azar, $P(L)=2/7$ y $P(L^c)=5/7$.

\textbf{a)} Por probabilidad total,
$P(V)=(2/7)(1/4)+(5/7)(2/5)=\boxed{5/14}$.

\textbf{b)} $P(L\cap V)=(2/7)(1/4)=1/14$.
Por Bayes, $P(L\mid V)=(1/14)/(5/14)=\boxed{1/5}$.

\textbf{c)} Entre los días sin lluvia, la probabilidad de que
tampoco haga viento es $1-0.40=0.60$. Por tanto
$P(L^c\cap V^c)=(5/7)(0.60)=\boxed{3/7}$.
